In theory, we would design a Turing Machine, working within some polynomial boundaries, and prove that this machine decides (solves) the language (decision problem)
or computes the function. However, this approach will take ages.

Instead of going through Turing Machines, we could use an informal but more expressive tool the so called \textbf{pseudocode}.
\begin{example}
    e.g. Suppose we want to show the following problem to be computable in polynomial time:
    \begin{center}
        \textit{given two strings $x,y \in \{0,1\}^*$ determine if $x$ contains an instance of $y$ (is $y$ a substring of $x$?)}
    \end{center} \vspace{7pt}
    A pseudocode solving the problem above is the following. \vspace{7pt}

    \begin{center}
        \includegraphics[width=0.35\textwidth]{img/img4.png}
    \end{center}

    The most amount of work is done inside the \textit{if condition}: by pseudocode we cannot anymore state that every operation has the same unary cost as
    Turing Machines. In this case the amount of work done is at least equal to the size of $y$.
\end{example}
How could we be sure that the algorithm shown works in polynomial time? We must satisfy some constraints about that, divided into:
\begin{enumerate}
    \item The input can be easily encoded as a binary string. Notice, we use pseudocode as a short-cut, but the formal proof should be done by the construction of a Turing Machine.
    \item The total number of instructions executed is polynomially bounded.
    \item All intermediate results are polynomially bounded in length.
    \item Each instruction takes polynomial time to be simulated.
\end{enumerate}
\begin{example}
    e.g. Defining the polynomial complexity of the previous algorithm. \vspace{7pt}

    \begin{enumerate}
        \item The total number of instructions executed is bounded by $O(|x|)$.
        \item The intermediate results obtained over $i$ cannot be greater than $O(|x|)$.
        \item Comparing two strings of length $|y|$ can be done in polynomial time in $|y|$, thus polynomial in $| \lfloor (x,y) \rfloor |$
    \end{enumerate}
\end{example}